\songtitle{\textchinese{三個夏夜} (San ge xia ye)}

\tag*{2026}\tag{house}\tag{mandarin-lyrics}\tag{tres-noches-de-verano}
% -> electronica.tex?
\begin{notes}
Chinese adaptation of \emph{Tres noches de verano} (page \pageref{sec:tres-noches-de-verano}). Retrieved from Dvigitt's Suno page (V6-MINI, 14 November 2026).
\end{notes}
\begin{style}
House, four-on-the-floor at about 124 BPM; filtered synth pad, kick and light hi-hats build with a snare roll, rising filter, vocal chops and melodic synth lead into heavy wobble-bass dubstep-style drops. Female Mandarin mezzo delivers a rhythmic count chant; seasonal sections rise into drops on their final lines, with a kick-free pad-and-vocal winter breakdown and a stripped outro.
\end{style}
\begin{lyrics}
\songpart{Intro (filtered synth pad, no drums, building)}
\begin{chinese}
夏天、秋天、冬天、春天
\end{chinese}

\songpart{Verse 1: Summer (kick enters, light hi-hats)}
\begin{chinese}
夏天的夜晚 在泳池邊\\
第二個夏夜 在泳池邊\\
第三個夏夜 在泳池邊\\
我的人生 什麼都沒改變

\songpart{Build-up (snare roll, rising filter, vocal chop repeating "三個夏夜")}
一個、兩個、三個
\end{chinese}

\songpart{Drop (heavy wobble bass, sparse vocal)}
\begin{chinese}
三個夏夜 三個夏夜 三個夏夜
\end{chinese}

\songpart{Verse 2: Autumn}
\begin{chinese}
秋天的白天 在湖畔\\
第二個秋日 在湖畔\\
第三個秋日 在湖畔\\
我的真愛 還沒有出現
\end{chinese}

\songpart{Build-up (added melodic synth lead)}
\begin{chinese}
一個、兩個、三個
\end{chinese}

\songpart{Drop (heavy wobble bass, sparse vocal)}
\begin{chinese}
三個夏夜 三個夏夜 三個夏夜
\end{chinese}

\songpart{Verse 3: Winter (breakdown: kick drops out, pads and vocal only)}
\begin{chinese}
冬天的午後 在我家\\
第二個冬日 在我家\\
第三個冬日 在我家\\
世界依舊 繞著我轉
\end{chinese}

\songpart{Verse 4: Spring (kick returns, full energy)}
\begin{chinese}
春天的早晨 在花園\\
第二個春晨 在花園\\
第三個春晨 在花園\\
我仍在等 你的愛
\end{chinese}

\annotation{Final Drop (everything at once, then a stripped outro)}\\
\begin{chinese}
等你的愛 等你的愛
\end{chinese}
\end{lyrics}
